\documentclass[10pt,a4paper]{amsart}
\usepackage[margin=19mm]{geometry}
\usepackage{amsmath,amssymb,mathtools}
\usepackage[scr=rsfs,cal=euler]{mathalfa}
\usepackage{xfrac}
\usepackage[unicode,hidelinks,bookmarks=false]{hyperref}
\newcommand{\ie}{i.e.\ }
\newcommand{\cf}{cf.\ }
\newcommand{\resp}{respectively\ }
\newcommand{\Subsection}{\subsection}
\providecommand{\nobreakdash}{\nobreak\hbox{-}}
\setlength{\emergencystretch}{2em}
\multlinegap0pt
\usepackage{bm}
\usepackage{bbm}
\usepackage{dsfont}
\usepackage{xspace}
\usepackage{cancel}
\usepackage{mathtools}
\usepackage{amsthm}
\makeatletter
\@ifundefined{theorem}{\newtheorem{theorem}{Theorem}}{}
\@ifundefined{assumption}{\newtheorem{assumption}[theorem]{Assumption}}{}
\@ifundefined{lemma}{\newtheorem{lemma}[theorem]{Lemma}}{}
\makeatother
\makeatletter
\newcommand{\R}{\mathbb{R}}
\renewcommand{\eqref}[1]{(\ref{#1})}
\expandafter\ifx\csname proof\endcsname\relax
\newenvironment{proof}[1][Proof]{\noindent\textit{#1.} }{\hfill \rule{0.5em}{0.5em}}
\fi
\makeatother
\title[Spreading speeds for prey--predator systems in a shifting en]{Spreading speeds for prey--predator systems in a shifting environment: a short proof}
\author{Zhucheng Jin}
\date{Source: arXiv:2606.14059v2 (CC BY 4.0)}
\begin{document}
\maketitle

\noindent\emph{Source and licence.} Spreading speeds for prey--predator systems in a shifting environment: a short proof. Version arXiv:2606.14059v2 (CC BY 4.0), distributed under the Creative Commons Attribution 4.0 International licence, as recorded in the arXiv licence metadata for this exact version and shown on the abstract page https://arxiv.org/abs/2606.14059v2. Full rights notice: licenses/arxiv-2606.14059-CC-BY-4.0.txt.\par\smallskip

\noindent\textit{Editorial scope.} Counted unit: the Lemma whose source label is \texttt{Lem-estimate} in the article (source0.tex lines 523--528, source label \texttt{Lem-estimate}), delivered together with the article's own complete original proof of it (source0.tex lines 529--573, one \texttt{proof} environment of 45 lines). No mathematics is written, repaired or re-derived: statement and proof are the article's own text line for line. Required context retained uncounted: Pb (lines 154--160); Ass-fg (lines 216--231); rrr (lines 226--228); Ass-bound (lines 241--247); Ass-IC (lines 250--259); ev (lines 282--284); Thm1 (lines 292--303); Lem-bound (lines 467--470); Lem-liou (lines 487--499). Every definition, display and label of those lines is retained unchanged. Omitted from this package and not redistributed: 1--153, 161--215, 229--240, 248--249, 260--281, 285--291, 304--466, 471--486, 500--522, 574--950. Nothing inside a retained excerpt is omitted or altered: every retained body is delivered exactly as the article's lines are written, so no omission receipt of any kind accompanies this package. Citations: the retained excerpts carry no citation command at all, so no bibliography entry of the article is transcribed and no reference is redistributed. Reference adaptation: none was needed and none was made; every reference inside the delivered unit resolves inside it, so no unresolved reference or citation is produced. Printed slips kept literal: no printed word, symbol, hypothesis or number of the counted statement or of its proof was altered. Layout and class: the article's own class is \texttt{article} with packages including amssymb, amsfonts, graphicx, amsmath, appendix, hyperref, mathrsfs, frcursive, geometry, color; the wrapper is the lane's pinned \texttt{amsart} editorial wrapper at 10pt on A4, which restates the theorem-like environments the retained text needs and adds the article's own macro definitions that the retained text needs (\texttt{\textbackslash R}, \texttt{\textbackslash eqref}), with extra wrapper packages (bm, bbm, dsfont, xspace, cancel, mathtools). The delivered numbering is the unit's own. Assets: no retained span contains a figure environment, a graphics-inclusion command, a PDF-page inclusion, a table or a TikZ picture; no image file is copied into this package, the package declares no asset dependency and no image file is redistributed with it. Licence: version arXiv:2606.14059v2 of this article is distributed under the Creative Commons Attribution 4.0 International licence, as recorded in the arXiv licence metadata for this exact version and shown on the abstract page https://arxiv.org/abs/2606.14059v2. A full rights notice is delivered at licenses/arxiv-2606.14059-CC-BY-4.0.txt. Characters: every retained excerpt is pure ASCII, so the delivered engine, XeLaTeX with its default Latin Modern text font, reports no missing character.
\begin{equation}\label{Pb}
\begin{cases}
\partial_t u= d \partial_{xx} u  +  u f(x-c_1t, u,v),  & t>0, x\in \R,\\
\partial_t v= \; \partial_{xx} v + vg( x-c_1 t,  u,  v ),  & t>0, x\in \R, \\
u(0, x)=u_0(x), \;\; v(0, x)= v_0(x), & x\in \R,
\end{cases}
\end{equation}

\begin{assumption}\label{Ass-fg}
Assume that $f: \R \times [0, \infty)\times  [0, \infty) \to \R$ and $g: \R \times [0, \infty)\times  [0, \infty) \to \R$   satisfy: 
\begin{enumerate}
\item[(i)]  For every $R>0$,  $z \mapsto f(z,u,v)$ and $z \mapsto g(z,u,v )$ are bounded and uniformly continuous on $\R$, uniformly for $(u,v) \in [0, R]^2$.  Moreover,   $f$ and  $g $ are Lipschitz continuous in $(u,v)\in [0,R]^2$,  uniformly  with respect to $z\in \R$;
\item [(ii)]  (\textbf{Predation type}) For each $u>0$, $z\in \R$, the  map  $v\mapsto f(z,u, v)$ is decreasing. For each $v>0$,  $z\in \R$, the map $u\mapsto g(z,u,v)$  is increasing.
\item[(iii)] (\textbf{Monostable}) $f(\cdot, 1,0) \equiv 0$ and $\inf_{z\in \R} f(z,u, 0)> 0$ for all $u\in [0, 1)$. 
\item [(iv)] (\textbf{Strong KPP}) For each $u\geq 0$,  the map  $v \mapsto g(z, u, v)$ is nonincreasing uniformly in $z\in \R$.  
\item [(v)] The limits $\lim_{z\to \pm \infty} g(z, 1,0 )$ exist  and 
$$\inf_{z\in \R} g(z, 1, 0) >0. $$
 Denote
\begin{equation}\label{rrr}
r^{\pm} := g(\pm \infty , 1, 0).
\end{equation}

\end{enumerate}
\end{assumption}

\begin{assumption} \label{Ass-bound}
Assume that there exist constants $\overline{V}>0$  and $\delta>0$ such that
$$  \sup_{z\in\R} g(z,1,\overline{V})\leq 0   \quad
\text{ and  } \quad
\inf_{\substack{z\in \R, \; 0\leq u\leq \delta}} f(z,u, \overline {V})> 0.
$$
\end{assumption}

\begin{assumption}\label{Ass-IC}
Let  $\overline{V}$ be chosen  as in Assumption \ref{Ass-bound}.  Let $u_0$ and $v_0$ be two bounded and continuous functions satisfying
$$
0<\inf_{x\in\R} u_0(x)\le u_0(x) \le 1,
\quad
0\le v_0(x)\le \overline{V},\quad
v_0\not\equiv 0,  \quad \text{ for all }   x\in \R. 
$$
and suppose that $v_0$ is compactly supported. 
\end{assumption}

\begin{equation}\label{ev}
\Lambda_1 := \Lambda_1(g)= \inf\left\{ \Lambda \in \R: \;\;  \exists \psi \in C^2_{loc} (\R), \psi>0,  \partial_{zz} \psi + g(z, 1,0) \psi \leq \Lambda \psi \text{ in } \R  \right\}.
\end{equation}

\begin{theorem}\label{Thm1}
Let Assumptions \ref{Ass-fg}, \ref{Ass-bound} and \ref{Ass-IC} be satisfied and $c_1>0$.  Let $(u,v)$ be the solution of  \eqref{Pb}.   Then the  spreading speed of $v$ is given by 
\begin{equation*}
c_v^*= \begin{cases}
2\sqrt{r^+}, & \text{ if  }  0<c_1 \leq 2 \sqrt{r^+}, \\
c_1,   & \text{ if }  2\sqrt{r^+} \leq c_1 \leq  2 \sqrt{\Lambda_1}, \\
\frac{c_1}{2} - \sqrt{\Lambda_1 - r^{-} }  + \frac{r^{-}}{\frac{c_1}{2} - \sqrt{\Lambda_1 - r^{-} }  },   \quad  & \text{ if  }   2 \sqrt{\Lambda_1}  \leq  c_1 \leq 2 (  \sqrt{r^-} + \sqrt{\Lambda_1 -r^{-}}),\\
2\sqrt{r^-},    & \text{ if  } c_1  \geq   2 (  \sqrt{r^-} + \sqrt{\Lambda_1 -r^{-}}),
\end{cases}
\end{equation*}
 where $r^{\pm}$ and $\Lambda_1$ are defined in \eqref{rrr} and \eqref{ev} respectively.   
\end{theorem}

\begin{lemma}\label{Lem-bound}
Let Assumptions \ref{Ass-fg}, \ref{Ass-bound} and \ref{Ass-IC} be satisfied. Denote $\beta:= \min\{  \inf_{x\in\R } u_0(x),   \delta \}$.  Then  
$$\beta \leq  u(t,x)\leq 1,   \;\; 0 \leq v(t,x) \leq   \overline{V}, \quad  \forall  t\geq 0, x\in \R. $$  
\end{lemma}

\begin{lemma}\label{Lem-liou}
Let $u$  be a bounded entire solution of
 \begin{equation*}
  \partial_t u= d \partial_{xx} u +  u \tilde{f}(x-c_1 t, u, 0), \quad  (t,x) \in \R^2.
 \end{equation*} 
where   $\tilde f(\cdot, u, 0)$  is a locally uniform limit of translations of  $f(\cdot, u, 0)$
and satisfies the monostable condition inherited from Assumption \ref{Ass-fg}\textit{(iii)}. 
If
\[
0<\beta\leq u(t,x)\leq 1 \quad \text{in } \mathbb R^2,
\]
then $u \equiv 1$.
\end{lemma}

\begin{lemma}\label{Lem-estimate}
Let the assumptions of  Theorem \ref{Thm1} be satisfied.   For each $\alpha>0$, there exist  constants  $M_{\alpha}>0$ and $T_{\alpha} >0$ such that 
\begin{equation*}
1-u(t,x) \leq \alpha+ M_{\alpha} v(t,x), \quad \forall t\geq T_{\alpha}, \forall x\in\R.
\end{equation*}
\end{lemma}

\begin{proof}
Argue by contradiction and suppose that there exist  sequences $\{x_n\}_{n\geq 1} \subset \R $,  $\{t_n\}_{n\geq 1} \subset \R$ with  $t_n \to +\infty$ as $n \to +\infty$,     and some constant  $\alpha_0>0 $  such that 
\begin{equation} \label{contra1}
1-u(t_n, x_n) > \alpha_0 + n  v(t_n, x_n), \;\;\;\; \forall n\geq 1.
\end{equation}
Set
\begin{equation*}
u_n(t,x) := u(t+t_n, x+x_n) \text{ and } v_n(t,x) := v(t+t_n, x+x_n). 
\end{equation*}
Note that $(u_n, v_n) (t,x)$ satisfies 
\begin{equation*}
\begin{cases}
\partial_t u_n= d \partial_{xx} u_n +   u_n f \big( x+x_n -c_1(t+t_n), u_n, v_n \big), & \forall t>-t_n, x\in \R, \\
\partial_t v_n =  \partial_{xx} v_n   + v_n  g\big(  x+x_n -c_1(t+t_n), u_n,   v_n  \big), & \forall t>-t_n, x\in \R.
\end{cases}
\end{equation*}
By parabolic estimates and Assumption \ref{Ass-fg}\textit{(i)}, up to a subsequence, we have 
\begin{align*}
\lim_{n\to +\infty} u_n(t,x) = \tilde{u}(t,x) \quad \text{ locally uniformly in }  (t,x) \in \R^2, \\
\lim_{n\to +\infty} v_n(t,x) = \tilde{v}(t,x) \quad  \text{ locally uniformly in }  (t,x) \in \R^2,
\end{align*}
and 
\begin{align*}
\lim_{n\to +\infty} f \big( x+x_n -c_1(t+t_n), u_n, v_n \big) = \tilde{f} \big( x-c_1t, \tilde{u}, \tilde{v} \big) \quad \text{ locally uniformly in }  (t,x) \in \R^2, \\
\lim_{n\to +\infty} g \big( x+x_n -c_1(t+t_n), u_n, v_n \big) = \tilde{g} \big( x-c_1t, \tilde{u}, \tilde{v} \big) \quad \text{ locally uniformly in }  (t,x) \in \R^2,
\end{align*}
The parabolic estimates ensure that    $(\tilde{u}, \tilde{v}) (t,x)$  satisfies  for all $(t,x) \in \R^2$,
\begin{equation*}
\begin{cases}
\partial_t \tilde{u}= d \partial_{xx} \tilde{u}  +  \tilde{u}  \tilde{f} \big( x-c_1t, \tilde{u}, \tilde{v} \big),\\
\partial_t \tilde{v}=  \partial_{xx} \tilde{v} + \tilde{v} \tilde{g} \big( x-c_1t, \tilde{u}, \tilde{v} \big).  \\
\end{cases}
\end{equation*}
Letting $n \to +\infty $ in  \eqref{contra1} and using the boundedness of $u$,  we  obtain
$\tilde{v}(0, 0) = 0$. The  strong maximum principle   applied to  the $\tilde{v}$-equation  implies that $\tilde{v}\equiv 0$.  Thus, $\tilde{u}$ satisfies 
\begin{equation*}
\partial_t \tilde{u}= d \partial_{xx} \tilde{u}  + \tilde{u} \tilde{f} \big( x-c_1t, \tilde{u}, 0 \big) ,     \quad \forall (t,x ) \in \R^2.
\end{equation*}
Moreover,  Lemma \ref{Lem-bound} gives $\tilde{u} \geq \beta >0$.   
Then  the Liouville-type  Lemma \ref{Lem-liou}  gives  $\tilde{u} \equiv 1$.     However, \eqref{contra1} yields 
$1- \tilde{u}(0, 0) \geq \alpha_0 >0$.
This contradiction completes  the proof. 


\end{proof}

\end{document}
